\section{Gestione delle modifiche dopo la baseline}

Questa sezione documenta le modifiche significative apportate alla baseline, riportandone una breve sintesi e il collegamento alle relative GitHub Issue e/o Pull Request. La cronologia completa resta tracciata nel repository GitHub.

\subsection*{Tipi di modifica}
\begin{itemize}
    \item \textbf{Req./UC/US/AC/MU Change:} modifica di un requisito, use case, user story, acceptance criteria o mockup.
    \item \textbf{Clarification:} rende esplicito un significato precedentemente ambiguo o incompleto.
\end{itemize}

\subsection*{Registro delle modifiche}

\begin{longtable}{p{1.6cm} p{1.7cm} p{4.5cm} p{2.2cm} p{3cm} p{2.5cm}}
\caption{Registro delle modifiche di TEPO successive alla baseline.}\\
\toprule
\textbf{ID} & \textbf{Tipo} & \textbf{Descrizione} & \textbf{Origine} & \textbf{Artefatti impattati} & \textbf{GitHub / Stato} \\
\midrule
\endfirsthead
\toprule
\textbf{ID} & \textbf{Tipo} & \textbf{Descrizione} & \textbf{Origine} & \textbf{Artefatti impattati} & \textbf{GitHub / Stato} \\
\midrule
\endhead

CH-TEPO-01 & Clarification & La soglia di prestazione dell'ottimizzazione è riferita a 30 opzioni candidate (non 26 come inizialmente indicato), in coerenza con l'estensione prevista del catalogo. & Team review & RNF\ref{rnf:prestazioni}; US\ref{us:ottimizza}; AC-US07-03 & \#12 --- DONE \\
CH-TEPO-02 & Req. Change & Aggiunta esplicita, fra i limiti dichiarati, dell'assenza di pianificazione multi-anno e di frontiera di Pareto. & Stakeholder (debriefing) & Sezione 1; RNF\ref{rnf:isolamento} & \#18 --- DONE \\
CH-TEPO-03 & Clarification & Chiarito che la componente AI può ricadere in modo automatico sul motore a regole qualora il provider esterno non sia configurato o non risponda. & Team review & RF\ref{rf:ai_proposta}; RNF\ref{rnf:ai_isolamento}; UC-09; AC-US12-02 & \#22 --- DONE \\
CH-TEPO-04 & MU Change & Ridisegno del pannello di proposta AI con evidenziazione degli interventi proposti come esclusi senza deselezione automatica. & User feedback interno & MU-Fig.5; US\ref{us:ai_appr}; AC-US13-01 & \#25 --- OPEN \\
CH-TEPO-05 & Req. Change & Introdotto requisito di snapshot immutabile delle revisioni dati (RNF\ref{rnf:snapshot}) per garantire la ricostruibilità dei run di audit. & Audit interno & RNF\ref{rnf:snapshot}; RF\ref{rf:storico}; UC-12 & \#28 --- OPEN \\
CH-TEPO-06 & Clarification & Chiarito che la pagina di sensibilità esegue uno sweep deterministico sul budget e non una simulazione Monte Carlo. & Team review & RF\ref{rf:sensibilita}; US\ref{us:sensibilita}; AC-US11-01 & \#31 --- DONE \\

\bottomrule
\end{longtable}

\textbf{Nota operativa.} Dopo la baseline, il D1 è un \emph{living artifact}: le modifiche significative vengono tracciate su GitHub e sottoposte a impact analysis. Gli artefatti successivi (D2, codice, test) sono indicati come impattati solo se già esistenti al momento della modifica. La cronologia completa è mantenuta nel repository, questa sezione ne riporta la sintesi.